\section{Grafy a~jejich reprezentace}\label{sec:grafy}

Grafy jsou základním stavebním kamenem mnoha úloh v~informatice. Zachycují objekty a~jejich vzájemné vztahy. Lze jimi reprezentovat třeba propojení zařízení v~počítačové síti,~schéma elektrického obvodu,~mapu města nebo vazby mezi atomy. Jejich matematická podstata je však velice jednoduchá -- \emph{vrcholy} propojené \emph{hranami}.

\begin{definition}[Graf]
    Grafem $G$ nazveme uspořádanou dvojici $(V,~E)$,~kde $V$ je množina \emph{vrcholů} (nebo také \emph{uzlů}) a~$E$ množina \emph{hran},~přičemž pokud
    \begin{itemize}
        \item $E\subseteq\set{\set{u,v}\mid u,v\in V,~u\neq v}$,~pak $G$ nazýváme \emph{neorientovaným} grafem;
        \item $E\subseteq\set{(u,v)\mid u,v\in V,~u\neq v}$,~pak $G$ nazýváme \emph{orientovaným} grafem.
    \end{itemize}
\end{definition}

\begin{figure}[ht]
    \centering
    \begin{subfigure}{6.5cm}
        \centering
        \begin{tikzpicture}[scale=0.9]
    \node[vertex] (a) at (0,0) {$1$};
    \node[vertex] (b) at (2,1.4) {$2$};
    \node[vertex] (c) at (4,0) {$3$};
    \node[vertex] (d) at (2,-1.4) {$4$};
    \draw[diredge] (a)--(b);
    \draw[diredge] (b)--(c);
    \draw[diredge] (c)--(d);
    \draw[diredge] (a)--(d);
    \draw[diredge,bend left=18] (d) to (b);
\end{tikzpicture}

        \caption{Příklad orientovaného grafu.}
        \label{fig:graf_priklad1}
    \end{subfigure}
    \qquad
    \begin{subfigure}{6.5cm}
        \centering
        \begin{tikzpicture}[scale=0.9]
    \foreach \i in {1,...,6}
        \node[vertex] (v\i) at ({90-60*(\i-1)}:1.7) {$\i$};
    \foreach \a/\b in {1/2,1/3,1/4,1/6,2/3,2/5,3/4,3/6,4/5,5/6}
        \draw[edge] (v\a)--(v\b);
\end{tikzpicture}

        \caption{Příklad neorientovaného grafu.}
        \label{fig:graf_priklad2}
    \end{subfigure}
\end{figure}

\subsection{Incidence, dosažitelnost a~vzdálenost}

Je-li \(e\in E\) hrana a~\(v\in V\) jeden z~jejích koncových vrcholů, říkáme, že \(e\) je s~vrcholem \(v\) \emph{incidentní}. V~neorientovaném grafu jsou dva vrcholy sousední, vede-li mezi nimi hrana. V~orientovaném grafu rozlišujeme směr: hrana \((u,v)\) vychází z~\(u\) a~vstupuje do \(v\).

Vrchol \(v\) je \emph{dosažitelný} z~vrcholu \(u\), jestliže existuje cesta z~\(u\) do \(v\). Délkou cesty v~neohodnoceném grafu rozumíme počet jejích hran. Vzdálenost
\[
    d(u,v)
\]
je délka nejkratší cesty z~\(u\) do \(v\); pokud taková cesta neexistuje, klademe \(d(u,v)=\infty\). V~neorientovaném grafu platí \(d(u,v)=d(v,u)\), zatímco v~orientovaném grafu se obě hodnoty mohou lišit a~jedna z~nich může být nekonečná.

Pojmy dosažitelnosti a~vzdálenosti připravují dvě různé otázky. Prohledávání grafu může pouze zjistit, které vrcholy jsou ze startu dosažitelné, nebo může navíc hledat nejkratší cesty. BFS zvládne obě úlohy v~neohodnoceném grafu; u~ohodnocených hran budeme později potřebovat Dijkstrův algoritmus.

\subsection{Stupeň vrcholu a~princip sudosti}

V~neorientovaném grafu definujeme \emph{stupeň vrcholu} \(v\) jako počet hran incidentních s~\(v\) a~značíme jej \(\deg(v)\). U~orientovaného grafu budeme v~souladu s~prezentacemi značit vstupní stupeň \(\deg^+(v)\), výstupní stupeň \(\deg^-(v)\) a~celkový stupeň
\[
    \deg(v)=\deg^+(v)+\deg^-(v).
\]
Vstupní stupeň počítá hrany končící ve \(v\), výstupní stupeň hrany, které z~\(v\) vycházejí.

\begin{proposition}[Princip sudosti]\label{prop:princip-sudosti}
    Pro každý konečný neorientovaný graf \(G=(V,E)\) platí
    \[
        \sum_{v\in V}\deg(v)=2\sizeof{E}.
    \]
\end{proposition}

\begin{proof}
    Při sčítání stupňů započítáme každou hranu právě dvakrát, jednou u~každého jejího koncového vrcholu. Celkový součet je proto dvojnásobkem počtu hran.
\end{proof}

Z~principu sudosti plyne, že počet vrcholů lichého stupně je sudý. Součet všech stupňů je totiž sudý a~součet lichého počtu lichých čísel by byl lichý. Tato jednoduchá podmínka bývá užitečná při kontrole, zda zadaná posloupnost stupňů vůbec může pocházet z~nějakého grafu.

\begin{remark}
    Za počátek teorie grafů se obvykle považuje Eulerovo řešení úlohy o~sedmi mostech města Královec z~roku 1736. Euler nahradil části města vrcholy a~mosty hranami a~ukázal, že existence požadované procházky závisí pouze na stupních vrcholů, nikoli na přesném tvaru mapy.
\end{remark}

\subsection{Reprezentace grafu}

Podstatnou otázkou u~grafů jsou způsoby jejich reprezentace,~neboť při programování jsou různé způsoby různě efektivní a~je dobré si mezi nimi umět vybrat. Pracujme tedy s~grafem $G=(V,E)$,~přičemž vrcholy si pro jednoduchost označíme přirozenými čísly $1,2,\dots,n$,~tzn. $V=\set{1,2,\dots,n}$.
\begin{itemize}
    \item \textbf{Matice sousednosti.} První způsob,~jak graf reprezentovat,~je pomocí matice. V~tomto případě graf $G$ popisuje matice $n\times n$,~kde na pozici $ij$ je $1$,~pokud z~vrcholu $i$ do vrcholu $j$ vede hrana,~jinak $0$.
    \begin{example}[Prásátkový graf]
        Orientovaný graf $G$ (viz obrázek \ref{fig:graf_prasatko}) a~jeho matice sousednosti $\mathbf{A}$.
        \begin{figure}[ht]
            \centering
            \begin{tikzpicture}[scale=1]

\node[vertex,fill=white] (V0)   at (3,0)        {$v_0$};
\node[vertex,fill=white] (V1)   at (0,0)        {$v_1$};
\node[vertex,fill=white] (V2)   at (-1,0.5)     {$v_2$};
\node[vertex,fill=white] (V3)   at (0,-2)       {$v_3$};
\node[vertex,fill=white] (V4)   at (-0.5,-3)    {$v_4$};
\node[vertex,fill=white] (V5)   at (0.5,-3) {$v_5$};
\node[vertex,fill=white] (V6)   at (3,-2) {$v_6$};
\node[vertex,fill=white] (V7)   at (2.5,-3) {$v_7$};
\node[vertex,fill=white] (V8)   at (3.5,-3) {$v_8$};
\node[vertex,fill=white] (V9)   at (4.5,-1) {$v_9$};

\draw[diredge] (V0) -- (V1);
\draw[diredge] (V0) -- (V9);
\draw[diredge] (V0) -- (V3);

\draw[diredge] (V1) -- (V3);
\draw[diredge] (V1) -- (V2);

\draw[diredge] (V3) -- (V4);

\draw[diredge] (V5) -- (V3);

\draw[diredge] (V6) -- (V0);
\draw[diredge] (V6) -- (V3);
\draw[diredge] (V6) -- (V7);

\draw[diredge] (V8) -- (V6);

\draw[diredge] (V9) -- (V6);

\end{tikzpicture}
            \caption{Prasátkový graf.}
            \label{fig:graf_prasatko}
        \end{figure}
        \[\mathbf{A}=\left(\begin{matrix}
            0 & 1 & 0 & 1 & 0 & 0 & 0 & 0 & 0 & 1\\
            0 & 0 & 1 & 1 & 0 & 0 & 0 & 0 & 0 & 0\\
            0 & 0 & 0 & 0 & 0 & 0 & 0 & 0 & 0 & 0\\
            0 & 0 & 0 & 0 & 1 & 0 & 0 & 0 & 0 & 0\\
            0 & 0 & 0 & 0 & 0 & 0 & 0 & 0 & 0 & 0\\
            0 & 0 & 0 & 1 & 0 & 0 & 0 & 0 & 0 & 0\\
            1 & 0 & 0 & 1 & 0 & 0 & 0 & 1 & 0 & 0\\
            0 & 0 & 0 & 0 & 0 & 0 & 0 & 0 & 0 & 0\\
            0 & 0 & 0 & 0 & 0 & 0 & 1 & 0 & 0 & 0\\
            0 & 0 & 0 & 0 & 0 & 0 & 1 & 0 & 0 & 0
        \end{matrix}\right)\]
    \end{example}
    \item \textbf{Matice incidence.} Řádky jsou indexovány vrcholy a~sloupce jednotlivými hranami. Na místě $ij$ je $1$,~pokud hrana vede \emph{do} vrcholu $i$,~nebo $-1$,~pokud z~něj vychází,~jinak $0$. Pro neorientované grafy píšeme vždy $1$,~pokud daná hrana vede \emph{do/z} vrcholu.
    \begin{example}[Úplný graf $K_4$]
        Graf $G$ má $4$ vrcholy a~celkem $6$ hran (viz obrázek \ref{fig:graf_k4}); jeho matici incidence označíme $\mathbf{I}$. Řádky představují (shora dolů) vrcholy $0,1,2,3$ a~sloupce (zleva doprava) jednotlivé hrany $e_1,\dots,e_6$.
        \begin{figure}[ht]
            \centering
            \begin{tikzpicture}
    \node[vertex] (v0) at (0,1.6) {$0$};
    \node[vertex] (v1) at (-1.6,0) {$1$};
    \node[vertex] (v2) at (0,-1.6) {$2$};
    \node[vertex] (v3) at (1.6,0) {$3$};
    \foreach \a/\b in {v0/v1,v0/v2,v0/v3,v1/v2,v1/v3,v2/v3}
        \draw[edge] (\a)--(\b);
\end{tikzpicture}

            \caption{Úplný graf $K_4$.}
            \label{fig:graf_k4}
        \end{figure}
        \[\mathbf{I}=\left(\begin{matrix}
            1 & 0 & 0 & 1 & 1 & 0\\
            1 & 1 & 0 & 0 & 0 & 1\\
            0 & 1 & 1 & 0 & 1 & 0\\
            0 & 0 & 1 & 1 & 0 & 1
        \end{matrix}\right)\]
    \end{example}
    \item \textbf{Seznam sousedů.} Jedná se pravděpodobně o~nejpoužívanější variantu reprezentace grafů. Pro každý vrchol $v\in V$ si jednoduše pamatujeme seznam $S[v]$ všech vrcholů,~do kterých přímo vede hrana z~vrcholu $v$.\footnote{Název \emph{seznam} zde odkazuje na použitou datovou strukturu. Z~matematického pohledu popisuje $S[v]$ množinu sousedů.}
    \begin{example}
        Graf $G$ s~vrcholy $1,\dots,4$ a~jejich seznamy sousedů $S$ (viz obrázek \ref{fig:graf_seznam_sousedu}).
        \begin{figure}[ht]
            \centering
            \begin{subfigure}{6cm}
                \centering
                \begin{tikzpicture}
    \node[vertex] (v1) at (0,1.6) {$1$};
    \node[vertex] (v2) at (-1.5,0) {$2$};
    \node[vertex] (v3) at (1.5,0) {$3$};
    \node[vertex] (v4) at (0,-1.6) {$4$};
    \draw[diredge] (v1)--(v2);
    \draw[diredge] (v1)--(v3);
    \draw[diredge] (v2)--(v3);
    \draw[diredge] (v2)--(v4);
    \draw[diredge] (v3)--(v4);
\end{tikzpicture}

            \end{subfigure}
            \qquad
            \begin{subfigure}{6cm}
                \centering
                \begin{align*}
                    S[1]=\set{2,3}&\qquad S[3]=\set{4}\\
                    S[2]=\set{3,4}&\qquad S[4]=\emptyset
                \end{align*}
                \vfill
            \end{subfigure}
            \caption{Graf $G$ a~seznamy sousedů pro jeho každý vrchol.}
            \label{fig:graf_seznam_sousedu}
        \end{figure}
    \end{example}
\end{itemize}

Rozdíly shrnuje tabulka~\ref{tab:reprezentace-grafu}. Uvedené časy předpokládají přímou maticovou reprezentaci a~obyčejný nesetříděný seznam sousedů. Jiná datová struktura může některou operaci urychlit za cenu větší paměti nebo složitější aktualizace.

\begin{table}[ht]
    \centering
    \renewcommand{\arraystretch}{1.2}
    \begin{tabular}{lccc}
        \hline
        Reprezentace & Paměť & Test hrany \((u,v)\) & Výpis sousedů \(u\) \\
        \hline
        matice sousednosti & \(\bigO(n^2)\) & \(\bigO(1)\) & \(\bigO(n)\) \\
        matice incidence & \(\bigO(nm)\) & \(\bigO(m)\) & \(\bigO(nm)\) \\
        seznamy sousedů & \(\bigO(n+m)\) & \(\bigO(\deg(u))\) & \(\bigO(\deg(u))\) \\
        \hline
    \end{tabular}
    \caption{Základní porovnání reprezentací grafu.}
    \label{tab:reprezentace-grafu}
\end{table}

Matice sousednosti je vhodná pro husté grafy a~časté dotazy na existenci konkrétní hrany. Seznamy sousedů jsou přirozenou volbou pro řídké grafy a~algoritmy, které postupně procházejí všechny hrany. Matice incidence se v~běžné implementaci algoritmů používá méně, ale je důležitá v~algebraickém popisu grafů a~síťových úlohách. U~neorientovaného grafu uloží seznamy sousedů každou hranu dvakrát, jednou u~každého koncového vrcholu; konstantní násobek však v~\(\mathcal O\)-notaci nehraje roli.
